\documentclass{article}
\usepackage{amsmath,amssymb,amsthm}
\IfFileExists{hyperref.sty}{%
  \IfFileExists{color.sty}{\usepackage{color}\definecolor{paperlinkblue}{rgb}{0.10,0.25,0.42}%
    \usepackage[colorlinks=true,linkcolor=paperlinkblue]{hyperref}}{\usepackage[hidelinks]{hyperref}}%
}{}
\ifdefined\hyperlink
  \newcommand{\paperanchor}[1]{\hypertarget{#1}{}}%
  \newcommand{\paperlink}[2]{\hyperlink{#1}{\begin{NoHyper}#2\end{NoHyper}}}%
\else
  \ifdefined\pdfextension
    \ifnum\outputmode>0
      \def\paperdest{\pdfextension dest}\def\paperstart{\pdfextension startlink}%
      \def\paperend{\pdfextension endlink}\def\papercolor{\pdfextension colorstack}%
    \fi
  \else\ifdefined\pdfstartlink
    \ifnum\pdfoutput>0
      \let\paperdest\pdfdest\let\paperstart\pdfstartlink\let\paperend\pdfendlink\let\papercolor\pdfcolorstack
    \fi
  \fi\fi
  \ifdefined\paperdest
    \newcommand{\paperanchor}[1]{\paperdest name{#1} xyz}%
    \newcommand{\paperlink}[2]{\leavevmode\paperstart attr{/Border [0 0 0]} goto name{#1}%
      \papercolor 0 push {0.10 0.25 0.42 rg 0.10 0.25 0.42 RG}#2\papercolor 0 pop\paperend}%
  \else\ifdefined\XeTeXversion
    \newcommand{\paperanchor}[1]{\special{pdf:dest (#1) [@thispage /XYZ null @ypos null]}}%
    \newcommand{\paperlink}[2]{\leavevmode\special{pdf:bann << /Subtype /Link /Border [0 0 0] /A << /S /GoTo /D (#1) >> >>}%
      \special{color push rgb 0.10 0.25 0.42}#2\special{color pop}\special{pdf:eann}}%
  \else
    \PackageError{lean-to-latex}{No PDF hyperlink backend available}{Use XeLaTeX, pdfLaTeX, LuaLaTeX, or install hyperref.}%
  \fi\fi
\fi
\newcommand{\paperlineanchor}[1]{\raisebox{10pt}[0pt][0pt]{\paperanchor{#1}}}
\DeclareRobustCommand{\paperresultname}[1]{\paperlineanchor{theorem:\arabic{theorem}}#1}
\newtheoremstyle{paperresult}{8pt}{5pt}{\normalfont\raggedright}{}{\bfseries}{.}{\newline}{}
\theoremstyle{paperresult}
\newtheorem{theorem}{\paperresultname{Theorem}}
\newcounter{paperpremise}
\newtheorem{lemma}[theorem]{\paperresultname{Lemma}}
\newtheorem{corollary}[theorem]{\paperresultname{Corollary}}
\newcounter{definingequation}
\newenvironment{proofstep}{\par\smallskip\ifdim\leftskip<3em\relax\advance\leftskip by 1em\relax\fi\noindent\ignorespaces}{\par\smallskip}
\setlength{\textwidth}{6.5in}
\setlength{\oddsidemargin}{0in}
\AddToHook{env/proof/begin}{\raggedright}
\setlength{\emergencystretch}{\linewidth}
\title{A recurrence and a factorial quotient for uniform Bell numbers}
\author{}
\date{}
\begin{document}
\maketitle
\section*{Definitions}
\paragraph{Bell.} \par\noindent\refstepcounter{definingequation}\label{defeq:1}\paperlineanchor{defeq:1}\textup{(\paperlink{defeq:1}{D1})} Let \(m\) be a multiset of elements of \(\mathbb{N}\). Then \(\operatorname{bell}(m) = \operatorname{multinomial}(\operatorname{to\,Finset}(m),\allowbreak{} k \mapsto k \cdot \operatorname{count}(k,\allowbreak{} m)) \cdot \prod_{k\in \operatorname{to\,Finset}(m) \setminus \{0\}}(\prod_{0\le j<\operatorname{count}(k,\allowbreak{} m)}(\binom{j \cdot k + k - 1}{k - 1}))\).

\paragraph{Uniform bell.} \par\noindent\refstepcounter{definingequation}\label{defeq:2}\paperlineanchor{defeq:2}\textup{(\paperlink{defeq:2}{D2})} Let \(m\) and \(n\) be natural numbers. Then \(\operatorname{uniform\,Bell}(m,\allowbreak{} n) = \operatorname{bell}(\operatorname{replicate}(m,\allowbreak{} n))\).

\subsection*{Prerequisites}
\begin{description}
\item[\paperlineanchor{premise:46696e7365742e70726f645f636f6e6772}\paperlink{premise:46696e7365742e70726f645f636f6e6772}{P1.}] \refstepcounter{paperpremise}\label{premise:46696e7365742e70726f645f636f6e6772}\textbf{Equality of products with equal factors.} Let \(\iota\) and \(M\) be types. Let \(s_{1}\) and \(s_{2}\) be finite subsets of \(\iota\). Fix a structure \(\mathit{inst}\) of type \(\operatorname{Comm\,Monoid}\left(M\right)\). Let \(f\) and \(g\) be functions of type \(\iota \to M\). Assume \(s_{1} = s_{2}\). Assume \(\forall x\, :\, \iota,\; x \in s_{2} \implies f\left(x\right) = g\left(x\right)\). Then \(\prod_{i\in s_{1}}\left(f(i)\right) = \prod_{i\in s_{2}}\left(g(i)\right)\).
\item[\paperlineanchor{premise:46696e7365742e70726f645f72616e67655f73756363}\paperlink{premise:46696e7365742e70726f645f72616e67655f73756363}{P2.}] \refstepcounter{paperpremise}\label{premise:46696e7365742e70726f645f72616e67655f73756363}\textbf{Product range successor.} Let \(M\) be a type. Fix a structure \(\mathit{inst}\) of type \(\operatorname{Comm\,Monoid}\left(M\right)\). Let \(f\) be a function of type \(\mathbb{N} \to M\). Let \(n\) be a natural number. Then \(\prod_{0\le x<n + 1}(f(x)) = \prod_{0\le x<n}(f(x)) \cdot f(n)\).
\item[\paperlineanchor{premise:4e61742e6469765f65715f6f665f65715f6d756c5f6c656674}\paperlink{premise:4e61742e6469765f65715f6f665f65715f6d756c5f6c656674}{P3.}] \refstepcounter{paperpremise}\label{premise:4e61742e6469765f65715f6f665f65715f6d756c5f6c656674}\textbf{Div equals of equals product left.} Let \(n\), \(m\), and \(k\) be natural numbers. Assume \(0 < n\). Assume \(m = k \cdot n\). Then \(m / n = k\).
\item[\paperlineanchor{premise:4e61742e666163746f7269616c5f706f73}\paperlink{premise:4e61742e666163746f7269616c5f706f73}{P4.}] \refstepcounter{paperpremise}\label{premise:4e61742e666163746f7269616c5f706f73}\textbf{Factorial positivity.} Let \(n\) be a natural number. Then \(0 < n!\).
\item[\paperlineanchor{premise:4e61742e6d756c5f706f73}\paperlink{premise:4e61742e6d756c5f706f73}{P5.}] \refstepcounter{paperpremise}\label{premise:4e61742e6d756c5f706f73}\textbf{Product positivity.} Let \(n\) and \(m\) be natural numbers. Assume \(n > 0\). Assume \(m > 0\). Then \(n \cdot m > 0\).
\item[\paperlineanchor{premise:4e61742e706f775f706f73}\paperlink{premise:4e61742e706f775f706f73}{P6.}] \refstepcounter{paperpremise}\label{premise:4e61742e706f775f706f73}\textbf{Power positivity.} Let \(a\) and \(n\) be natural numbers. If \(0 < a\), then \(0 < a^{n}\).
\item[\paperlineanchor{premise:4e61742e756e69666f726d42656c6c5f6571}\paperlink{premise:4e61742e756e69666f726d42656c6c5f6571}{P7.}] \refstepcounter{paperpremise}\label{premise:4e61742e756e69666f726d42656c6c5f6571}\textbf{Uniform bell equals.} Let \(m\) and \(n\) be natural numbers. Then \(\operatorname{uniform\,Bell}\left(m, n\right) = \prod_{0\le p<m}\left(\binom{p \cdot n + n - 1}{n - 1}\right)\).
\item[\paperlineanchor{premise:4e61742e756e69666f726d42656c6c5f6d756c5f6571}\paperlink{premise:4e61742e756e69666f726d42656c6c5f6d756c5f6571}{P8.}] \refstepcounter{paperpremise}\label{premise:4e61742e756e69666f726d42656c6c5f6d756c5f6571}\textbf{Uniform bell product equals.} Let \(m\) and \(n\) be natural numbers. Assume \(n \ne 0\). Then \(\operatorname{uniform\,Bell}\left(m, n\right) \cdot n!^{m} \cdot m! = \left(m \cdot n\right)!\).
\item[\paperlineanchor{premise:65715f636f6d6d}\paperlink{premise:65715f636f6d6d}{P9.}] \refstepcounter{paperpremise}\label{premise:65715f636f6d6d}\textbf{Equals commutation.} Let \(\alpha\) be a sort. Let \(a\) and \(b\) be elements of \(\alpha\). Then \(a = b \iff b = a\).
\item[\paperlineanchor{premise:6d756c5f6173736f63}\paperlink{premise:6d756c5f6173736f63}{P10.}] \refstepcounter{paperpremise}\label{premise:6d756c5f6173736f63}\textbf{Multiplication associativity.} Let \(G\) be a type. Fix a structure \(\mathit{inst}\) of type \(\operatorname{Semigroup}\left(G\right)\). Let \(a\), \(b\), and \(c\) be elements of \(G\). Then \(a \cdot b \cdot c = a \cdot \left(b \cdot c\right)\).
\item[\paperlineanchor{premise:6d756c5f636f6d6d}\paperlink{premise:6d756c5f636f6d6d}{P11.}] \refstepcounter{paperpremise}\label{premise:6d756c5f636f6d6d}\textbf{Multiplication commutativity.} Let \(G\) be a type. Fix a structure \(\mathit{inst}\) of type \(\operatorname{Comm\,Magma}\left(G\right)\). Let \(a\) and \(b\) be elements of \(G\). Then \(a \cdot b = b \cdot a\).
\end{description}

\section{Results}
\begin{theorem}[{Recurrence for one additional block}]\label{theorem:1}
Let \(m\) and \(n\) be natural numbers. Then
\begin{center}
\(\displaystyle \operatorname{uniform\,Bell}(m + 1,\allowbreak{} n) = \binom{m \cdot n + n - 1}{n - 1} \cdot \operatorname{uniform\,Bell}(m,\allowbreak{} n).\)
\end{center}
\end{theorem}

\begin{proof}
The substitution uses the following argument.
\begin{proofstep}
Fix \(x\) in \(\mathbb{N}\).
Assume \(x \in \{i\in\mathbb{N}\mid i<m + 1\}\).
By \paperlink{premise:6d756c5f636f6d6d}{Prerequisite~P\ref{premise:6d756c5f636f6d6d}} with \(G = \mathbb{N}\) and \(x\) in place of its parameter \(a\) and \(b = n\), we obtain \(x \cdot n = n \cdot x\).
Applying the same function to equal arguments, we obtain \(\binom{x \cdot n + n - 1}{n - 1} = \binom{n \cdot x + n - 1}{n - 1}\).
\end{proofstep}
The substitution uses the following argument.
\begin{proofstep}
Fix \(x\) in \(\mathbb{N}\).
Assume \(x \in \{i\in\mathbb{N}\mid i<m\}\).
By \paperlink{premise:6d756c5f636f6d6d}{Prerequisite~P\ref{premise:6d756c5f636f6d6d}} with \(G = \mathbb{N}\) and \(x\) in place of its parameter \(a\) and \(b = n\), we obtain \(x \cdot n = n \cdot x\).
Applying the same function to equal arguments, we obtain \(\binom{x \cdot n + n - 1}{n - 1} = \binom{n \cdot x + n - 1}{n - 1}\).
\end{proofstep}
Rewriting with \paperlink{premise:4e61742e756e69666f726d42656c6c5f6571}{Prerequisite~P\ref{premise:4e61742e756e69666f726d42656c6c5f6571}}, \paperlink{premise:46696e7365742e70726f645f636f6e6772}{Prerequisite~P\ref{premise:46696e7365742e70726f645f636f6e6772}}, \paperlink{premise:46696e7365742e70726f645f72616e67655f73756363}{Prerequisite~P\ref{premise:46696e7365742e70726f645f72616e67655f73756363}} (the final factor at index \(m\) is \(\binom{n \cdot m + n - 1}{n - 1}\)), \paperlink{premise:6d756c5f636f6d6d}{Prerequisite~P\ref{premise:6d756c5f636f6d6d}} and the two identities established above gives \(\operatorname{uniform\,Bell}(m + 1,\allowbreak{} n) = \binom{m \cdot n + n - 1}{n - 1} \cdot \operatorname{uniform\,Bell}(m,\allowbreak{} n)\).
\end{proof}

\begin{theorem}[{Factorial quotient for uniform Bell numbers}]\label{theorem:2}
Let \(m\) and \(n\) be natural numbers. Assume \(n \ne 0\). Then
\begin{center}
\(\displaystyle \operatorname{uniform\,Bell}(m,\allowbreak{} n) = (m \cdot n)! / (n!^{m} \cdot m!).\)
\end{center}
\end{theorem}

\begin{proof}
By \paperlink{premise:4e61742e666163746f7269616c5f706f73}{Prerequisite~P\ref{premise:4e61742e666163746f7269616c5f706f73}}, we obtain \(0 < n!\).
By \paperlink{premise:4e61742e706f775f706f73}{Prerequisite~P\ref{premise:4e61742e706f775f706f73}}, we obtain \(0 < n!^{m}\).
By \paperlink{premise:4e61742e666163746f7269616c5f706f73}{Prerequisite~P\ref{premise:4e61742e666163746f7269616c5f706f73}} with \(m\) in place of its parameter \(n\), we obtain \(0 < m!\).
By \paperlink{premise:4e61742e6d756c5f706f73}{Prerequisite~P\ref{premise:4e61742e6d756c5f706f73}}, we obtain \(n!^{m} \cdot m! > 0\).
By definition, we obtain \(\operatorname{uniform\,Bell}(m,\allowbreak{} n) \cdot n!^{m} \cdot m! = \operatorname{uniform\,Bell}(m,\allowbreak{} n) \cdot n!^{m} \cdot m!\).
Rewriting with \paperlink{premise:4e61742e756e69666f726d42656c6c5f6d756c5f6571}{Prerequisite~P\ref{premise:4e61742e756e69666f726d42656c6c5f6d756c5f6571}} (using \(n \ne 0\)) in the preceding result gives \(\left(m \cdot n\right)! = \operatorname{uniform\,Bell}\left(m, n\right) \cdot n!^{m} \cdot m!\).
Rewriting with \paperlink{premise:6d756c5f6173736f63}{Prerequisite~P\ref{premise:6d756c5f6173736f63}} in the preceding result gives \(\left(m \cdot n\right)! = \operatorname{uniform\,Bell}\left(m, n\right) \cdot \left(n!^{m} \cdot m!\right)\).
By \paperlink{premise:4e61742e6469765f65715f6f665f65715f6d756c5f6c656674}{Prerequisite~P\ref{premise:4e61742e6469765f65715f6f665f65715f6d756c5f6c656674}}, we obtain \(\left(m \cdot n\right)! / \left(n!^{m} \cdot m!\right) = \operatorname{uniform\,Bell}\left(m, n\right)\).
Rewriting with \paperlink{premise:65715f636f6d6d}{Prerequisite~P\ref{premise:65715f636f6d6d}} in the preceding result gives \(\operatorname{uniform\,Bell}\left(m, n\right) = \left(m \cdot n\right)! / \left(n!^{m} \cdot m!\right)\).
\end{proof}

\end{document}
